\section{实验原理}

\subsection{文件系统布局}
文件系统存放在磁盘中,多数磁盘划分为一个或多个分区,每个分区有一个独立的文件系统。磁盘分区和文件系统中的基本组成如下:
\begin{itemize}
    \item 主引导记录(Master Boot Record,MBR):位于磁盘的0号扇区,用来引导计算机,其后面是分区表,该表给出每个分区的起始和结束地址。
    \item 引导块(boot block):MBR执行引导块中的程序后,该程序负责启动该分区中的操作系统。
    \item 超级块(super block):包含文件系统的所有关键信息,在计算机启动时会被读入内存。超级块中包含分区的快的数量、块的大小、空闲分区的数量和指针、空闲FCB数量和FCB指针。
    \item 文件系统空闲块信息:可以使用位示图或指针链接的形式给出;后面跟着一组i节点(每个文件对应一个节点);接着是根目录,存放在文件系统目录树的根部;最后是其他文件所有目录和文件。
\end{itemize}

\subsection{FAT}
显式链接是文件分配物理块的一种方式(类似静态链表),其中将链接文件所有块的指针
显示地存入内存的一张表中,即文件分配表(FAT,File Allocation Table),每个表项
存放链接指针,即下一盘块号。因此只要将文件的第一个盘块号存放在FCB中,后续盘块号都可以通过查FAT表来获取。
因此FAT的表项与全部磁盘块一一对应,可以使用一个特殊数字如$-2$表示空闲块,用$-1$表示文件结束块。
也就是说FAT不仅记录了文件各块之间的前后链接关系,还标记了空闲块信息,还可以用来实现文件管理如分配空闲磁盘块等。


FAT 的优点是：FAT 表将文件各数据块之间的链接关系集中保存在一张表中，并且通常会在系统启动时被读入内存。
因此，系统在查找文件后续数据块时，只需要在内存中查询 FAT 表，而不需要反复访问磁盘中的前驱数据块，从而减少磁盘访问次数，提高文件检索效率。
与此同时，FAT 表还能够统一记录空闲块信息，便于系统完成磁盘空间的分配与回收。

FAT 的缺点是：由于 FAT 表需要记录全部磁盘块的状态，因此磁盘容量较大时，FAT 表本身会占用较多内存空间。
此外，虽然 FAT 表能够减少磁盘访问次数，但当系统只知道文件的起始块号时，要访问文件中较靠后的数据块，仍然需要沿 FAT 链逐项查找。
因此，FAT 的随机访问效率仍然低于连续分配方式。


\subsection{文件的打开与关闭}
当用户对文件进行操作时,每次都需要从检索目录开始,这将大大降低文件操作的效率。因此文件在使用前需要通过系统调用open被显式地打开。操作系统维护一张包含所有文件信息的表(打开文件表)。
这里"打开"的含义是指调用open根据文件名搜索目录,将指明文件的属性从外存复制到内存打开文件表的一个表目中,并将该表目的编号(索引,也即文件描述符)返回给用户。当用户再次向系统发出文件操作请求时,
通过索引在打开文件表中查到文件信息,从而节省再次搜索目录的开销。当文件不再使用时,利用系统调用close关闭文件,实质即为操作系统从打开文件表中删除该文件的条目。

操作系统中存在多个进程可以同时打开文件时,采用两级表:整个系统表和每个进程表。整个系统的打开文件表包含FCB的副本级其他信息,每个进程的打开文件表包含指向系统表中
适当条目的指针。因此系统打开文件表为每个文件关联一个打开计数器(Open Count),以记录有多少个进程打开了该文件。每个关闭操作close 使count递减,当打开计数器为0时
表示该文件不再被使用,就可从系统打开文件表中删除相应条目。